\section{Conclusion}

SecFlowOps executes real scanner, remediation, policy, metric, adjudication, DAST, and performance workflows. The full controlled protocol executes 432/432 planned local runs across three repository families, eight scenarios, six configurations, and three repetitions. The controlled size study, the 75-run open-source campaign, the NodeGoat npm-remediation study, and the injected external recall study provide complementary checks on size effects, external repositories, package-manager remediation, and injected ground-truth matching. The ZAP probe validates executable OWASP ZAP scanning, JSON normalization, and reflected-XSS ground-truth matching under the Java 24 runtime available in the artifact environment.

The resulting contribution is an evaluable DevSecOps framework and measurement protocol, not a production-ready autonomous remediation system. The evidence supports claims about local scanner execution, controlled and injected recall, partial external remediation, policy gating, ZAP-enabled DAST execution, measured overhead, remote GitHub Actions execution, real pull-request creation, owner merge acceptance, measured merge latency, and post-merge CI. It does not support complete recall over naturally occurring OSS vulnerabilities, exploitability of every scanner finding, compatibility preservation for all remediated applications, independent human review quality, or production deployment outside the repository.
